\documentclass[10pt,a4paper]{amsart}
\usepackage[margin=19mm]{geometry}
\usepackage{amsmath,amssymb,mathtools}
\usepackage[scr=rsfs,cal=euler]{mathalfa}
\usepackage{xfrac}
\usepackage[unicode,hidelinks,bookmarks=false]{hyperref}
\newcommand{\ie}{i.e.\ }
\newcommand{\cf}{cf.\ }
\newcommand{\resp}{respectively\ }
\newcommand{\Subsection}{\subsection}
\providecommand{\nobreakdash}{\nobreak\hbox{-}}
\setlength{\emergencystretch}{2em}
\multlinegap0pt
\usepackage{bm}
\usepackage{bbm}
\usepackage{dsfont}
\usepackage{xspace}
\usepackage{cancel}
\usepackage{mathtools}
\usepackage{amsthm}
\makeatletter
\@ifundefined{theorem}{\newtheorem{theorem}{Theorem}}{}
\@ifundefined{proposition}{\newtheorem{proposition}[theorem]{Proposition}}{}
\makeatother
\title[Automorphic side of Taylor-Wiles method for orthogonal and s]{Automorphic side of Taylor-Wiles method for orthogonal and symplectic groups}
\author{Xiaoyu Zhang}
\date{Source: arXiv:2411.04897v2 (CC BY 4.0)}
\begin{document}
\maketitle

\noindent\emph{Source and licence.} Automorphic side of Taylor-Wiles method for orthogonal and symplectic groups. Version arXiv:2411.04897v2 (CC BY 4.0), distributed under the Creative Commons Attribution 4.0 International licence, as recorded in the arXiv licence metadata for this exact version and shown on the abstract page https://arxiv.org/abs/2411.04897v2. Full rights notice: licenses/arxiv-2411.04897-CC-BY-4.0.txt.\par\smallskip

\noindent\textit{Editorial scope.} Counted unit: the Proposition whose source label is \texttt{None} in the article (source0.tex lines 2321--2345, source label \texttt{None}), delivered together with the article's own complete original proof of it (source0.tex lines 2346--2780, one \texttt{proof} environment of 435 lines). No mathematics is written, repaired or re-derived: statement and proof are the article's own text line for line. Required context retained uncounted: none. Every definition, display and label of those lines is retained unchanged. Omitted from this package and not redistributed: 1--2320, 2781--11750. Nothing inside a retained excerpt is omitted or altered: every retained body is delivered exactly as the article's lines are written, so no omission receipt of any kind accompanies this package. Citations: the retained excerpts carry no citation command at all, so no bibliography entry of the article is transcribed and no reference is redistributed. Reference adaptation: none was needed and none was made; every reference inside the delivered unit resolves inside it, so no unresolved reference or citation is produced. Printed slips kept literal: no printed word, symbol, hypothesis or number of the counted statement or of its proof was altered. Layout and class: the article's own class is \texttt{amsart} with packages including geometry, inputenc, tgtermes, babel, amsmath, mathtools, amssymb, graphicx, bbm, xcolor, amsthm, mathrsfs, imakeidx, mathdots; the wrapper is the lane's pinned \texttt{amsart} editorial wrapper at 10pt on A4, which restates the theorem-like environments the retained text needs and adds the article's own macro definitions that the retained text needs (none), with extra wrapper packages (bm, bbm, dsfont, xspace, cancel, mathtools). The delivered numbering is the unit's own. Assets: no retained span contains a figure environment, a graphics-inclusion command, a PDF-page inclusion, a table or a TikZ picture; no image file is copied into this package, the package declares no asset dependency and no image file is redistributed with it. Licence: version arXiv:2411.04897v2 of this article is distributed under the Creative Commons Attribution 4.0 International licence, as recorded in the arXiv licence metadata for this exact version and shown on the abstract page https://arxiv.org/abs/2411.04897v2. A full rights notice is delivered at licenses/arxiv-2411.04897-CC-BY-4.0.txt. Characters: every retained excerpt is pure ASCII, so the delivered engine, XeLaTeX with its default Latin Modern text font, reports no missing character.
	\begin{proposition}
		There is a finite set of elements
		$(b_i)_{i\in I}$ in $G(F_v)$
		depending on $j$ and $n_0$,
		independent of $m$,
		such that
		\[
		U_0^{(n_0)}(v^m)
		\varsigma_{v,j}
		U_0^{(n_0)}(v^m)
		=
		\bigsqcup_{i\in I}
		b_i
		U_0^{(n_0)}(v^m),
		\]
		\[
		U_1^{(n_0)}(v^m)
		\varsigma_{v,j}
		U_1^{(n_0)}(v^m)
		=
		\bigsqcup_{i\in I}
		b_i
		U_1^{(n_0)}(v^m).
		\]
	\end{proposition}

	\begin{proof}
		To ease notations,
		we write
		$U_?$
		for $U_?^{(n_0)}(v^m)$
		($?=0,1$)
		and
		$\varsigma=\varsigma_{v,j}$.
		We have the following decomposition of
		$U_0$:
		write
		$U_{0,0}$
		for the subgroup of
		$U_0$
		consisting of those $g$ such that
		$g
		\equiv
		\mathrm{diag}
		(A,B,C)(\mathrm{mod}\,\varpi_v^m)$
		with $A,C\in
		B_{n_0}
		(\mathcal{O}_{F,v}/\varpi_v^m)$,
		$U_{0,1}$
		for the subgroup of
		$U_0$
		consisting of those 
		$g$ such that
		$g
		=
		\begin{pmatrix}
			1_{n_0} & \ast & \ast \\
			0 & 1_{n-2n_0} & \ast \\
			0 & 0 & 1_{n_0}
		\end{pmatrix}
		(\mathrm{mod}\,\varpi_v^m)$
		is unipotent.
		Then we have
		a semi-direct product
		\[
		U_0
		=
		U_{0,1}\rtimes U_{0,0}.
		\]
		Similarly, for the case of
		$U_1$,
		we write
		$U_{1,0}
		=
		U_{0,0}\cap
		U_1$
		and
		$U_{1,1}
		=
		U_{0,1}\cap
		U_1$.
		Note that $U_{1,1}=U_{0,1}$.
		Then we also have a semi-direct product
		\[
		U_1
		=
		U_{1,1}\rtimes U_{1,0}.
		\]
		To decompose the double coset
		$U_?
		\varsigma
		U_?$
		into right cosets,
		it suffices to choose representatives in
		$U_?$
		of the quotient
		$U_?/
		\big(
		U_?
		\cap
		\varsigma
		U_?\varsigma^{-1}
		\big)$.
		Therefore to prove the lemma,
		it suffices to show that
		there is a common set of representatives
		$(u_i)_{i\in I}$ inside
		$U_1$
		such that
		$U_?\varsigma U_?=
		\sqcup_iu_i\varsigma U_?$
		($?=0,1$).
		For $?=0,1$,
		we have
		$\varsigma
		U_?
		\varsigma^{-1}
		=
		\varsigma U_{?,1}\varsigma^{-1}
		\rtimes
		\varsigma U_{?,0}\varsigma^{-1}$,
		and then
		$U_?\cap
		\varsigma U_?\varsigma^{-1}
		=
		(
		U_{?,1}\cap
		\varsigma
		U_{?,1}\varsigma^{-1}
		)
		\rtimes
		(
		U_{?,0}\cap
		\varsigma
		U_{?,0}
		\varsigma^{-1}
		)$,
		from which one deduces
		\[
		U_?/
		(U_?\cap
		\varsigma U_?\varsigma^{-1})
		=
		U_{?,1}/
		(U_{?,1}\cap
		\varsigma
		U_{?,1}\varsigma^{-1})
		\cdot
		U_{?,0}/
		(U_{?,0}\cap
		\varsigma
		U_{?,0}
		\varsigma^{-1}).
		\]
		Since
		$U_{1,1}=U_{0,1}$,
		it suffices to show that there is a common set of
		representatives
		$(v_k)_{k\in J}$
		inside
		$U_{1,0}$
		such that
		$U_{?,0}\varsigma U_{?,0}
		=
		\sqcup_kv_k\varsigma U_{?,0}$
		($?=0,1$).
		Write
		$\varsigma=
		\mathrm{diag}(a_1,a_2,a_3)$
		with
		$a_1,a_3\in
		B_{n_0}
		(F_v)$.
		We also write an element
		$b\in
		U_{0,0}$
		in the form of block matrix
		$b=(b_{i,j})_{i,j=1}^3$
		with
		$b_{1,1},b_{3,3}$ of size
		$n_0\times n_0$
		and
		$b_{2,2}$ of size
		$(n-2n_0)\times(n-2n_0)$.
		There are two cases to consider
		\begin{enumerate}
			\item 
			If $j\geq n_0$
			(thus $a_1=1=a_3$),
			then for $?=0,1$,
			$U_{?,0}\cap
			\varsigma U_{?,}\varsigma^{-1}$
			consists of those
			$b=(b_{i,j})_{i,j=1}^3\in U_{?,0}$ with
			\[
			\begin{cases*}
				b_{k,l},\,
				a_kb_{k,l}a_l^{-1} 
				\equiv
				0
				(\mathrm{mod}\,\varpi_v^m),
				&
				$k\neq l$,
				\\
				b_{k,k},\,
				a_kb_{k,k}a_k^{-1}
				(\mathrm{mod}\,\varpi_v^m)
				\in
				\mathrm{GL}_{n-2n_0}(\mathcal{O}_{F,v}/\varpi_v^m),
				&
				$k=2$,
				\\
				b_{k,k}
				(\mathrm{mod}\,\varpi_v^m)
				\begin{cases*}
					\in
					B_{n_0}
					(\mathcal{O}_{F,v}/\varpi_v^m),
					&
					$?=0$,
					\\
					=1,
					&
					$?=1$,
				\end{cases*}
				&
				$k=1,3$.
			\end{cases*}
			\]
			Note that elements $b$ in
			$U_{1,0}$
			are determined by the blocks
			$b_{k,l}$
			with
			$k<l$.
			Set
			$r_1=n_0,r_2=n-2n_0,r_3=n_0$.
			For $1\leq k<l\leq3$,
			we choose a set of representatives
			$I_{k,l}$ in
			$\mathrm{M}_{r_k,r_l}
			(\mathcal{O}_{F,v})$
			of the quotient
			$\mathrm{M}_{r_k,r_l}
			(\mathcal{O}_{F,v})/
			(\mathrm{M}_{r_k,r_l}\cap
			a_k\mathrm{M}_{r_k,r_l}(\mathcal{O}_{F,v})a_l^{-1})$.
			Then the set
			\begin{equation}\label{set of representatives-1}
				I=
				\big\{
				b\in U_{1,0}\mid 
				b_{1,1}=1=b_{3,3},
				b_{k,l}\in I_{k,l},
				\forall k<l,
				\,
				\&
				b_{k,l}=0,
				\forall
				k>l
				\big\}
			\end{equation}
			is a set of representatives for both
			$U_{1,0}/(U_{1,0}\cap
			\varsigma U_{1,0}\varsigma^{-1})$	
			and
			$U_{0,0}/(U_{0,0}\cap
			\varsigma U_{0,0}\varsigma^{-1})$:
			indeed, for the first case,
			note that
			for any
			$b\in
			U_{1,0}$,
			the matrices
			$b_{k,l}
			(\mathrm{mod}\,\varpi_v^m)$
			($k<l$)
			is necessarily the image of an element
			$b_{k,l}'$ in
			$I_{k,l}$
			mod $\varpi_v^m$.
			Now take an element
			$b'\in I$
			with
			$b'_{k,l}=b_{k,l}$
			for $k<l$
			(the existence is guaranteed by the surjectivity of
			$G(\mathcal{O}_{F,v})
			\rightarrow
			G(\mathcal{O}_{F,v}/\varpi_v^m)$),
			then it is clear that
			$(b')^{-1}b\in
			U_{1,0}\cap
			\varsigma U_{1,0}\varsigma^{-1}$;
			for the second case,
			for any $b\in U_{0,0}$,
			one can find an element
			$c=\mathrm{diag}(c_{1,1},c_{2,2},c_{3,3})
			\in
			U_{0,0}$
			such that
			$c_{k,k}
			\equiv
			b_{k,k}
			(\mathrm{mod}\,\varpi_v^m)$
			for all $k$.
			We set
			$d=bc^{-1}
			\in
			U_{1,0}$.
			As above we take
			$b'\in I$
			such that
			$(b')_{k,l}\equiv
			d_{k,l}
			(\mathrm{mod}\,\varpi_v^m)$
			for $k<l$.
			Then
			$(b')^{-1}d\in
			U_{1,0}
			\cap
			\varsigma
			U_{1,0}\varsigma^{-1}$
			and therefore
			$(b')^{-1}b\in
			U_{0,0}\cap
			\varsigma U_{0,0}\varsigma^{-1}$.
			Note that
			$I$ is independent of $m$.
			
			
			
			
			\item 
			If
			$j<n_0$,
			the argument is similar.
			For $k=1,3$,
			write
			$a_k=\mathrm{diag}((a_k)_1,(a_k)_2)$ in block matrix
			with
			$(a_1)_1,(a_3)_2$ of size $j\times j$.
			Similarly, write
			$b_{k,k}=((b_{k,k})_{s,t})_{s,t=1,2}$
			with
			$(b_{1,1})_{1,1},
			(b_{3,3})_{2,2}$ of size
			$j\times j$.
			For
			$?=0,1$,
			$U_{?,0}\cap
			\varsigma
			U_{?,0}\varsigma^{-1}$
			consists of those $b\in U_{?,0}$ with
			\[
			\begin{cases*}
				b_{k,l},\,
				a_kb_{k,l}a_l^{-1}		
				(\mathrm{mod}\,\varpi_v^m)
				\equiv
				0,
				&
				$k\neq l$,
				\\
				b_{k,k},\,
				a_kb_{k,k}a_k^{-1}
				(\mathrm{mod}\,\varpi_v^m)
				\in
				\mathrm{GL}_{n-2n_0}(\mathcal{O}_{F,v}/\varpi_v^m),
				&
				$k=2$,
				\\
				(b_{1,1})_{1,1},\,
				(b_{3,3})_{2,2}
				(\mathrm{mod}\,\varpi_v^m)
				\begin{cases*}
					\in
					B_j
					(\mathcal{O}_{F,v}/\varpi_v^m),
					&
					$?=0$,
					\\
					\equiv1,
					&
					$?=1$,
				\end{cases*}
				&
				\\
				(b_{1,1})_{2,2},\,
				(b_{3,3})_{1,1}
				(\mathrm{mod}\,\varpi_v^m)
				\begin{cases*}
					\in
					B_{n_0-j}
					(\mathcal{O}_{F,v}/\varpi_v^m),
					&
					$?=0$,
					\\
					\equiv1,
					&
					$?=1$,
				\end{cases*}
				\\
				(b_{1,1})_{2,1},\,
				(b_{3,3})_{1,2}^\mathrm{t}
				(\mathrm{mod}\,\varpi_v^{m+1})
				\equiv
				0.
				&
			\end{cases*}
			\]
			Note that
			$b_{3,3}
			(\mathrm{mod}\,\varpi_v^m)$
			is determined by
			$b_{1,1}
			(\mathrm{mod}\,\varpi_v^m)$.
			We fix a set of representatives
			$(I_{1,1})_{1,2}$
			in
			$\mathrm{M}_{j,n_0-j}
			(\mathcal{O}_{F,v})$
			for the quotient
			$\mathrm{M}_{j,n_0-j}
			(\mathcal{O}_{F,v}/\varpi_v)$.
			Write
			$I_{1,1}$
			for the set of elements
			$b_{1,1}\in
			\mathrm{M}_{j,j}(\mathcal{O}_{F,v})$
			such that
			$(b_{1,1})_{1,1}=(b_{1,1})_{2,2}=1,
			(b_{1,1})_{1,2}\in (I_{1,1})_{1,2},
			(b_{1,1})_{2,1}=0$.
			Then we take
			\begin{equation}\label{set of representatives-2}
				I:=
				\{
				b\in
				U_{1,0}\mid 
				b_{1,1}\in I_{1,1},
				b_{k,l}\in I_{k,l},
				\forall
				k<l,
				\,
				\&
				b_{k,l}=0,
				\forall
				k>l
				\}.
			\end{equation}
			Now one can verify as in the preceding case that
			$I$ is a common set of representatives for the quotients
			$U_{?,0}/
			(U_{?,0}\cap
			\varsigma U_{?,0}\varsigma^{-1})$
			($?=0,1$).
			Again we see that $I$ is independent of
			$m$.
		\end{enumerate}
	\end{proof}

\end{document}
